\section{The two first degeneracy correspondences}
The first degeneracy correspondence is already governed by relation (IV).
Put $B=\operatorname{PSL}_2(\mathbb Z)$ and write
\[
 u=\begin{pmatrix}1&1\\0&1\end{pmatrix},\qquad
 v=\begin{pmatrix}1&0\\-1&1\end{pmatrix}
\]
for their classes in $B$. The first map from $\Gamma_0(N)$ to $B$ is inclusion.
The second is conjugation by $D_N=\operatorname{diag}(N,1)$; denote it by $d_N$.
Thus $d_N(u)=u^N$ and $d_N(v^N)=v$. These maps extend to full profinite completions,
since each of their images has finite index in $B$.

\begin{proposition}\label{prop:degeneracy-two}
Let $F=(1,f)$ be a normalized element of $\widehat{\operatorname{GT}}$.
Let $\Phi_F$ be its standard action on $\widehat B$, so that
\[
 \Phi_F(u)=u,\qquad
 \Phi_F(v)=q^{-1}vq,\qquad q=f(u^2,v^2).
\]
Then $\Phi_F$ preserves the closure of $\Gamma_0(2)$, and its restriction
commutes with both degeneracy maps to $\widehat B$.
\end{proposition}
\begin{proof}
The standard index-three subgroup $\Gamma_0(2)$ is generated by $u,v^2$;
its presentation is $\langle u,w\mid(uw)^2=1\rangle$ with $w=v^2$.
This follows, for example, by applying the Schreier method to the action of
$B$ on $\mathbb P^1(\mathbb F_2)$: the stabilizer of infinity is $\Gamma_0(2)$,
and the two resulting generators can be chosen as $u,v^2$.
Because $q$ belongs to $\overline{\langle u^2,v^2\rangle}$, the formulas for
$\Phi_F$ take both generators into $\overline{\Gamma_0(2)}$.
The same argument for $F^{-1}$ gives equality of this subgroup with its image.
Commutation with inclusion is now its definition. For the other map,
\[
 d_2(q)=f(u^4,v)=f(u^2,v^2)=q.
\]
The middle equality follows by applying the interchange of $u,v$ to normalized
relation (IV), then applying $f(b,a)=f(a,b)^{-1}$ to both sides.
On the generator $u$ the two compositions both have value $u^2$.
On $v^2$ they both have value $q^{-1}vq$. Equality on these topological
generators proves the assertion.
\end{proof}

\section{Normalized congruence triviality and the level-three gate}
\label{sec:normalized-congruence}

The level-two identity has a consequence at every congruence level.
Use the signed representation
$\widetilde q:\widehat B_3\to\operatorname{SL}_2(\widehat{\mathbb Z})$
with $u\mapsto U=\left(\begin{smallmatrix}1&1\\0&1\end{smallmatrix}\right)$
and $v\mapsto V=\left(\begin{smallmatrix}1&0\\-1&1\end{smallmatrix}\right)$.

\begin{proposition}\label{prop:normalized-congruence-triviality}
Let $g=(1,f)\in\widehat{\mathrm{GT}}$ satisfy $\rho_2(g)=0$.
Then
\[
 f(U^2,V^2)=I
       \qquad\text{in }\operatorname{SL}_2(\widehat{\mathbb Z}).
\]
Consequently $\widetilde q\Phi_g=\widetilde q$ on $\widehat B_3$.
\end{proposition}
\begin{proof}
Normalized relation (IV), together with interchange of the two braid
generators and the two-cycle relation, gives
\[
 f(u^4,v)=f(u^2,v^2)
       \qquad\text{in }\widehat B_3.
\]
Fix a prime $\ell$ and put $Q=f(U^2,V^2)$ in
$\operatorname{SL}_2(\mathbb Z_\ell)$.
Conjugation by $D=\operatorname{diag}(2,1)$ in
$\operatorname{GL}_2(\mathbb Q_\ell)$ sends $U^2$ to $U^4$ and
$V^2$ to $V$. Therefore the displayed signed identity gives
$DQD^{-1}=Q$, with no scalar ambiguity. Since the two diagonal
entries of $D$ are distinct, $Q=\operatorname{diag}(a,a^{-1})$
for a unit $a\in\mathbb Z_\ell^\times$.

The images of the braid generators under
$\widetilde q_\ell\Phi_g$ are $U$ and
\[
 M=Q^{-1}VQ=\begin{pmatrix}1&0\\-a^2&1\end{pmatrix}.
\]
Their braid relation $UMU=MUM$, in its lower-left entry, is
$\gamma(1+\gamma)=0$ with $\gamma=-a^2$.
The coefficient ring is a domain and $a$ is a unit, so $a^2=1$.
Thus $Q=\pm I$ and $M=V$.

It remains to determine the sign using the three-cycle relation.
Put $A=U^2$, $B=V^2$, $Z=(AB)^{-1}$, and $R=UV$.
Direct multiplication gives
\[
 RAR^{-1}=B,\qquad RBR^{-1}=-Z,\qquad RZR^{-1}=-A.
\]
A commutator word is unchanged when one of its arguments is multiplied
by a central scalar. Since $Q$ is central, conjugating its defining
expression by $R$ and $R^2$ consequently gives
\[
 f(A,B)=f(B,Z)=f(Z,A)=Q.
\]
The multiplier-one three-cycle relation, with $ABZ=1$, now implies
$Q^3=I$. Together with $Q^2=I$, this forces $Q=I$.
This holds for every $\ell$, proving the signed adelic assertion.
Both braid generators now have unchanged images, which proves the
intertwining identity.
\end{proof}

In the following statements the ambient modular group is the
centerless group $B=\operatorname{PSL}_2(\mathbb Z)$ and its full
profinite completion, rather than the central braid extension. Put
\[
 \mathcal C=\ker\bigl(\widehat B\longrightarrow
      \operatorname{SL}_2(\widehat{\mathbb Z})/\{\pm I\}\bigr),
\]
where the scalar sign is the common sign in all components.

\begin{lemma}\label{lem:parabolic-centralizer}
The centralizer of $v$ in $\widehat B$ is
$\overline{\langle v\rangle}\simeq\widehat{\mathbb Z}$, and
it has trivial intersection with $\mathcal C$.
\end{lemma}
\begin{proof}
The free pure subgroup is
$\widehat F_2=\overline{\langle x\rangle}\amalg
\overline{\langle y\rangle}$, where $x=u^2,y=v^2$.
The centralizer theorem for a nonidentity element of a factor in a
free profinite product gives
$C_{\widehat F_2}(y)=\overline{\langle y\rangle}$
\cite[Theorem B, p.~155; Theorem B$'$, \S4, p.~159]{HerfortRibes1985}.
This theorem has no torsion assumption on the element of the factor.
If $c\in C_{\widehat B}(v)$, its image in $S_3$ centralizes the
transposition represented by $v$. Thus $cv^{-e}$ is pure for
some $e\in\{0,1\}$ and centralizes $y$. It lies in
$\overline{\langle y\rangle}$, proving the centralizer assertion.
Finally $v^t$ has congruence image represented by
$\left(\begin{smallmatrix}1&0\\-t&1\end{smallmatrix}\right)$.
It is a scalar sign only if $t=0$. This proves injectivity of the
procyclic subgroup in the congruence quotient and the last assertion.
\end{proof}

\begin{corollary}\label{cor:level-three-exact-gate}
For a normalized $g=(1,f)\in\widehat{\mathrm{GT}}$, put
$q=f(u^2,v^2)\in\widehat B$. The element $q$ lies in
$\mathcal C$, hence in $\overline{\Gamma_0(3)}$.
The automorphism $\Phi_g$ preserves that subgroup. For
$d_3:\overline{\Gamma_0(3)}\to\widehat B$ induced by
conjugation by $\operatorname{diag}(3,1)$, set
\[
 h_3=d_3(q)q^{-1}.
\]
Then $h_3\in\mathcal C$, and strict compatibility with both
level-three degeneracy maps is equivalent to the exact identity
\[
 d_3(q)=q.
\]
\end{corollary}
\begin{proof}
The preceding proposition shows that $q$ and the action of $g$
are trivial at every congruence level. Therefore $q$ belongs to
$\mathcal C$ and $\Phi_g$ preserves the level-three subgroup.
On its congruence image, $d_3$ is conjugation by the same diagonal
matrix: the equality holds on the discrete subgroup and extends
by continuity. Consequently $d_3(q)$ also has trivial congruence
image, and $h_3\in\mathcal C$.

The subgroup $\Gamma_0(3)$ is generated by $u,v^3$, and
$d_3(u)=u^3$, $d_3(v^3)=v$. The two proposed compositions agree
on $u$. On $v^3$ they agree exactly when
$h_3\in C_{\widehat B}(v)$. The preceding lemma and
$h_3\in\mathcal C$ make this equivalent to $h_3=1$.
\end{proof}

The corollary identifies the remaining full profinite identity. It does
not prove it: every congruence representation already kills $h_3$, so
congruence matrices alone cannot decide whether this element is trivial
in $\widehat B$.


\section{One local Brauer class for the two frame signs}
The two frame signs admit one local Brauer invariant which can be defined
from the traces of the two matrices. No choices of centralizer generators
or common matrix basis enter its definition.

\begin{proposition}\label{prop:local-brauer-frame}
Let
\[
 \mathcal D=\{(T,W)\in\operatorname{SL}_2(\mathbb Z_3)^2:
       \operatorname{tr}(T)\operatorname{tr}(W^2)\in\mathbb Z_3^\times\}.
\]
On this domain define
\[
 \mathfrak b(T,W)=
 \big[(\operatorname{tr}(T)\operatorname{tr}(W^2),3)_{\mathbb Q_3}\big]
       \in\operatorname{Br}(\mathbb Q_3)[2].
\]
Here brackets denote a quaternion algebra's Brauer class.
This invariant is unchanged by independent conjugation of $T$ and $W$.
Suppose their reductions have the form
\[
 \bar T\in(-I)^{\epsilon_3}\langle R_0\rangle,
 \qquad \bar W\in S_0^{\epsilon_2}\langle-I\rangle,
\]
where $R_0$ is the standard nonidentity unipotent matrix of order three
in $\operatorname{SL}_2(\mathbb F_3)$ and $S_0^2=-I$.
Then $(T,W)\in\mathcal D$ and
\[
 \mathfrak b(T,W)=(\epsilon_2+\epsilon_3)
                 [(-1,-3)_{\mathbb Q_3}].
\]
Its local Hilbert sign is $(-1)^{\epsilon_2+\epsilon_3}$.
\end{proposition}
\begin{proof}
Every power of $R_0$ has trace $2$ in $\mathbb F_3$, including its
zeroth power. Also $\bar W^2=(-I)^{\epsilon_2}$. Consequently
\[
 \operatorname{tr}(T)\equiv2(-1)^{\epsilon_3},\qquad
 \operatorname{tr}(W^2)\equiv2(-1)^{\epsilon_2}\pmod3.
\]
Both traces are units and their product has residue
$(-1)^{\epsilon_2+\epsilon_3}$. By Hensel's lemma, a unit of
$\mathbb Z_3$ is a square in $\mathbb Q_3$ precisely when its residue
is a square in $\mathbb F_3$. The odd-prime Hilbert formula gives
$(e,3)_3=(e/3)$ for a unit $e$. In particular $(-1,-3)_3=-1$,
so this quaternion algebra represents the unique nonzero class of
$\operatorname{Br}(\mathbb Q_3)[2]$. The formula follows.
Both traces are unchanged by independent conjugation, which proves
its invariance on the entire stated domain.
\end{proof}

For comparison with centralizer coordinates, suppose
\[
 T=aI+bR,\quad W=cI+dS,\qquad
 R^2+R+I=0,\quad S^2=-I,
 \qquad a^2-ab+b^2=c^2+d^2=1,
\]
where $R,S$ are the standard cubic and harmonic matrices over $\mathbb Z_3$.
The trace identities give
\[
 \operatorname{tr}(T)=2a-b\equiv2(a+b),\qquad
 \operatorname{tr}(W^2)=2(c^2-d^2)\pmod3.
\]
The norm equations imply that $a+b$ and $c^2-d^2$ are units. Thus
\[
 \mathfrak b(T,W)=\big[((a+b)(c^2-d^2),3)_{\mathbb Q_3}\big].
\]
This recovers the product of the two previously given local symbols,
while the trace expression makes the definition intrinsic to the ordered
pair of matrices.

The connection with the modular frame character is exact on its stated
domain. For $F\in H_0$, take the signed three-adic images of $A_F$ and
$\beta_F$ as $T$ and $W$. Their reductions satisfy the displayed conditions
with $\epsilon_3=a_3(F)$ and $\epsilon_2=a_2(F)$, by the definition of the
frame character. Therefore
\[
 \mathfrak b(T,W)=\delta(F)[(-1,-3)_{\mathbb Q_3}].
\]
Here $H_0=\{F:\lambda_F=1,\ t_F=0\pmod{12}\}$; this conclusion uses only
the finite frame conditions and does not require the uncorrected
three-adic matrices themselves to centralize the elliptic matrices.
For corrected frame pairs satisfying the centralizer norm forms above,
the same formula holds with their specified signed frame bits.

The class can be nonzero on the matrix domain: $T=-I,W=I$ gives its
nonzero value. This matrix calculation alone does not prove that such a
pair comes from any designated subgroup. On normalized elements of
$\widehat{\operatorname{GT}}$, the class vanishes by \cite[Theorem~3]{H17}. The algebra $(-1,-3)$ over $\mathbb Q$
has ramification set $\{3,\infty\}$, whereas $(2,3)$ has ramification set
$\{2,3\}$; these rational algebras cannot be identified merely because
their completions at three are isomorphic.

\section{The quartic Kummer layer over the Gaussian field}
The mod-four Kummer extension is a classical dihedral extension.
The following computation fixes its Hecke conductors and specifies the
additional analytic hypothesis that a square-root error term would use.
Let $K=\mathbb Q(i)$, let $\mathfrak p=(1+i)$, and put
$L=K(\alpha)$ with $\alpha=2^{1/4}>0$.
For a prime ideal $\mathfrak q\nmid2$ of $K$, use arithmetic Frobenius and define
\[
 \psi(\mathfrak q)=\frac{\operatorname{Frob}_{\mathfrak q}(\alpha)}{\alpha}
  =\left(\frac{2}{\mathfrak q}\right)_4.
\]
The last expression is the unique fourth root of unity whose reduction modulo
$\mathfrak q$ is $2^{(N\mathfrak q-1)/4}$. Extend multiplicatively to ideals
prime to two.

\begin{proposition}\label{prop:quartic-conductor}
The extension $L/K$ is cyclic of degree four. Its two faithful finite-order
Hecke characters are $\psi,\psi^{-1}$ and
\[
 \mathfrak f(\psi)=\mathfrak f(\psi^{-1})=\mathfrak p^6,
 \qquad \mathfrak f(\psi^2)=\mathfrak p^4.
\]
All three characters have trivial infinity type. The absolute discriminant of
$L$ is $2^{24}$. The representation
$\operatorname{Ind}_{G_K}^{G_{\mathbb Q}}\psi$ is odd and irreducible of
dimension two, has conductor $256$, and has determinant $\chi_{-4}$.
\end{proposition}
\begin{proof}
The polynomial $X^4-2$ is Eisenstein at two. Its real root field has degree four
and does not contain $i$, so $[L:\mathbb Q]=8$ and $[L:K]=4$.
The automorphism $r(\alpha)=i\alpha$, $r(i)=i$ generates
$\operatorname{Gal}(L/K)$; complex conjugation $s$ satisfies $srs=r^{-1}$.
This proves the dihedral description. Artin reciprocity identifies the
characters of this finite cyclic quotient with finite-order Hecke characters.
The displayed Frobenius formula follows by reducing $\alpha^{N\mathfrak q-1}$.
There is no ramification away from two, since the polynomial discriminant
of $X^4-2$ is a power of two and $K/\mathbb Q$ is unramified away from two.

We compute the different at two directly. Let $M=\mathbb Q(\zeta_8)=K(\sqrt2)$,
let $\varpi=1-\zeta_8$, and normalize the local valuation on $M$ by
$v_M(\varpi)=1$. Then $v_M(2)=4$ and $v_M(\sqrt2)=2$.
Put $D=\sqrt2/\varpi^2$. From
\[
 \sqrt2=\zeta_8+\zeta_8^{-1}
        =\frac{2-2\varpi+\varpi^2}{1-\varpi}
\]
we obtain $D\equiv1+\varpi\pmod{\varpi^2}$. For
$\beta=\alpha/\varpi$ and $\gamma=\beta-1$ we have
\[
 \gamma^2+2\gamma+(1-D)=0.
\]
This polynomial is Eisenstein over $M$ because $v_M(1-D)=1$.
Thus $L/M$ is totally ramified quadratic at two, $\gamma$ is a uniformizer,
and its different is generated by the derivative $2\gamma+2=2\beta$.
Since $\beta$ is a unit and $v_L(2)=8$, the different exponent is eight.
The discriminant of $M/\mathbb Q$ is $2^8$, from the cyclotomic polynomial
$X^4+1$. The discriminant formula in a tower therefore gives
$|d_L|=(2^8)^2\,2^8=2^{24}$.

Since $d_K=-4$, the relative discriminant ideals are
\[
 \mathfrak d_{M/K}=\mathfrak p^4,
 \qquad \mathfrak d_{L/K}=\mathfrak p^{16}.
\]
Indeed their norms are respectively $2^8/4^2=2^4$ and
$2^{24}/4^4=2^{16}$. The conductor-discriminant formula
\cite[Chapter V, Theorem~3.27, p.~165]{MilneCFT}
gives the conductor of the quadratic character $\psi^2$ as $\mathfrak p^4$.
Writing the common conductor exponent of the inverse faithful characters as
$a$, the same formula gives $2a+4=16$, hence $a=6$.

Conjugation by $s$ sends $\psi$ to $\psi^{-1}\ne\psi$, so induction is
irreducible. On $r$ its two eigenvalues are $i,-i$, and on $s$ they are
$1,-1$. Its determinant is consequently the character of $K/\mathbb Q$,
and it is odd. The induction formula for conductors gives
$|d_K|N_{K/\mathbb Q}\mathfrak f(\psi)=4\cdot2^6=256$.
\end{proof}

Hecke's theta construction associates to $\psi$ the weight-one form
\[
 \theta_\psi(z)=\sum_{\substack{\mathfrak a\subset\mathbb Z[i]\\
                     (\mathfrak a,2)=1}}
                 \psi(\mathfrak a)e^{2\pi i N\mathfrak a z}.
\]
It is the classical newform of level $256$ and character $\chi_{-4}$, recorded
as LMFDB $256.1.c.a$ \cite{LMFDB256}. Its first coefficients are
$q+q^9-2q^{17}-q^{25}+O(q^{33})$.

For an unramified prime ideal $\mathfrak q$ of $K$, specifying the value
$\psi(\mathfrak q)=i^a$ is a cyclic Chebotarev condition. A square-root
error term obtained from the explicit formula uses the named hypothesis
that every nontrivial zero of $L_K(s,\psi^j)$, $j=0,1,2,3$, has real part
$1/2$, where $j=0$ means $\zeta_K(s)$. The factorization
\[
 L_K(s,\psi^2)=L_{\mathbb Q}(s,\chi_8)
                    L_{\mathbb Q}(s,\chi_{-8}),\qquad
 \zeta_K(s)=\zeta_{\mathbb Q}(s)L_{\mathbb Q}(s,\chi_{-4})
\]
shows the new analytic input at this layer to be the quartic Hecke
$L$-function (the inverse character has the same $L$-function).
If extra cyclotomic congruences are imposed, the corresponding finite family
of Hecke twists $\psi^j(\chi\circ N_{K/\mathbb Q})$ must be included.
The existence of appropriate primes, and a terminating search for one, does
not require GRH. Nor does the algebraic decomposition itself require it.
A quantitative assertion must state its requested bound before a stronger
analytic hypothesis is described as necessary.
